\documentclass[11pt,letterpaper,reqno]{amsart}
\usepackage{tikz}
\usetikzlibrary{positioning, shapes.geometric, arrows.meta, calc, positioning}
\usepackage{amssymb}
\usepackage{amsmath}
\usepackage{amsthm}
\usepackage{amsfonts}
\usepackage{bbm}
\usepackage{enumitem} 
\usepackage{pgfplots}
\pgfplotsset{compat=1.18} 
\usepackage{booktabs}
\usepackage{graphicx}
\usepackage[T1]{fontenc}
\usepackage{doi}
\usepackage{float} 
\usepackage{comment} 
\addtolength{\hoffset}{-1.5cm}\addtolength{\textwidth}{3cm}
\addtolength{\voffset}{-1cm}\addtolength{\textheight}{2cm}

\usepackage{bookmark}
\usepackage{hyperref}
\hypersetup{pdfstartview={FitH}}
\newcommand{\C}{\mathbb{C}}
\newcommand{\cE}{\mathcal{E}}
\newcommand{\norm}[1]{\lVert #1 \rVert}
\newcommand{\abs}[1]{| #1 |}
\newcommand{\bv}{\mathbf{v}}
\newcommand{\bw}{\mathbf{w}}
\newcommand{\tr}{\operatorname{Tr}}
\DeclareMathOperator{\rank}{rank}




\newtheorem{theorem}{Theorem}[section]
\newtheorem{lemma}[theorem]{Lemma}
\newtheorem{proposition}[theorem]{Proposition}
\newtheorem{corollary}[theorem]{Corollary}
\newtheorem{claim}{Claim}
\newtheorem{question}[theorem]{Question}
\newtheorem{problem}[theorem]{Problem}
\newtheorem{conjecture}[theorem]{Conjecture}
\theoremstyle{definition}
\newtheorem{example}[theorem]{Example}
\newtheorem{remark}[theorem]{Remark}
\newtheorem{definition}[theorem]{Definition}
\numberwithin{equation}{section}

\newcommand{\snorm}[1]{\lVert#1\rVert}
\newcommand{\sang}[1]{\langle #1 \rangle}
\newcommand{\pvec}[1]{{\vec{#1}\,}'}
\newcommand{\avec}[1]{{\vec{#1}\,}^\ast}
\newcommand{\evec}[2]{{\vec{#1}\,}^{#2}}

\newcommand{\bs}{\boldsymbol}
\newcommand{\mb}{\mathbb}
\newcommand{\mbf}{\mathbf}
\newcommand{\mbm}{\mathbbm}
\newcommand{\mc}{\mathcal}
\newcommand{\mf}{\mathfrak}
\newcommand{\mr}{\mathrm}
\newcommand{\msc}{\mathscr}
\newcommand{\ol}{\overline}
\newcommand{\on}{\operatorname}
\newcommand{\wh}{\widehat}
\newcommand{\wt}{\widetilde}
\newcommand{\edit}[1]{{\color{red}#1}}
\newcommand{\imod}[1]{~\mathrm{mod}~#1}


% \left(\right) should behave the same as ()
\let\originalleft\left
\let\originalright\right
\renewcommand{\left}{\mathopen{}\mathclose\bgroup\originalleft}
\renewcommand{\right}{\aftergroup\egroup\originalright}
\newcommand{\tO}{\wt{\Omega}}
\newcommand{\wtilde}{\widetilde}
\renewcommand{\Re}{\mathrm{Re}}

\newcommand{\NN}{\mathbb{N}}
\newcommand{\taufunc}{\tau}
\newcommand{\omegap}{\omega}
\newcommand{\ord}{\operatorname{ord}}
\newcommand{\R}{\mathbb{R}}        % real numbers
\newcommand{\E}{\mathbb{E}}        % expectation
\newcommand{\Var}{\mathrm{Var}}    % variance
\newcommand{\Cov}{\operatorname{Cov}}
\newcommand{\PP}{\mathbb{P}}     % probability
\newcommand{\eps}{\varepsilon}     % epsilon
\newcommand{\ind}{\mathbf{1}}      % indicator function
\newcommand{\seq}[1]{\left(#1\right)} % sequence
\newcommand{\lcm}{\operatorname{lcm}}
\newcommand{\li}{\operatorname{li}}
\newcommand{\leg}[2]{\left(\frac{#1}{#2}\right)}
\newcommand{\A}{\mathcal{A}}
\newcommand{\muo}{\mu}
\newcommand{\ph}{\varphi}
\newcommand{\ta}{\tau}

\makeatother


\begin{document}


\title{A note on Problem~\#650}





\date{\today}



\maketitle









\section{Introduction}



In~\cite{Er95c}, Erd\H{o}s and Graham posed the following question, which also appears as Problem~\#650 on Bloom's Erd\H{o}s Problems website~\cite{EP}.
\begin{problem}\label{prob:650}
Let $f(m)$ be such that if $A\subseteq \{1,\ldots,N\}$ has $\lvert A\rvert=m$ then every interval in $[1,\infty)$ of length $2N$ contains $\geq f(m)$ many distinct integers $b_1,\ldots,b_r$ where each $b_i$ is divisible by some $a_i\in A$, where $a_1,\ldots,a_r$ are distinct.

Estimate $f(m)$. In particular is it true that $f(m)\ll m^{1/2}$?
\end{problem}



Erd\H{o}s and Sar\'{a}nyi \cite{ErSa59} proved that $f(m)\gg m^{1/2}$.




In this note, we solve Problem~\ref{prob:650}.






\section{Main Result}



\begin{theorem}
For every integer $m \ge 1$, one has
\[
f(m)\le 2\lceil \sqrt m\rceil.
\]
More generally, for all integers $s,t\ge 1$,
\[
f(st)\le s+t.
\]
In particular,
\[
f(m)\ll \sqrt m.
\]
\end{theorem}

\begin{proof}
We first prove the stronger statement
\[
f(st)\le s+t
\qquad (s,t\ge 1).
\]


Here, by a matching we mean a collection of distinct pairs $(a_i,b_i)$ such that
\[
a_i\in A,\qquad b_i\in I\cap \mathbb Z,\qquad a_i\mid b_i,
\]
with all the $a_i$ distinct and all the $b_i$ distinct.

If $s=1$ or $t=1$, then any matching has size at most $|A|=st$, so trivially
\[
f(st)\le st\le s+t.
\]
Thus it remains to consider the case $s,t\ge 2$.

Define
\[
E:=\prod_{p\le st} p^{\,1+\lfloor \log_p(s-1)\rfloor},
\qquad
D:=E\cdot \mathrm{lcm}(1,2,\dots,st).
\]
Then for every prime $p\le st$ and every integer $\delta$ with
\[
1\le |\delta|\le s-1,
\]
we have
\[
v_p(D)>v_p(\delta),
\]
where $v_p(\cdot)$ denotes the $p$-adic valuation. Indeed,
\[
v_p(\delta)\le \lfloor \log_p(s-1)\rfloor,
\]
while the exponent of $p$ in $E$ is
\[
1+\lfloor \log_p(s-1)\rfloor.
\]

Next let
\[
\mathcal P:=
\Bigl\{
p>st:\ p \text{ is prime and } p\mid (qD+\delta)
\text{ for some }1\le q\le t-1,\ 1\le |\delta|\le s-1
\Bigr\}.
\]
This is a finite set of primes.

For each $p\in\mathcal P$, consider the residue classes
\[
i-jD \pmod p
\qquad (1\le i\le s,\ 0\le j\le t-1).
\]
There are at most $st$ such classes, whereas $p>st$, so they do not exhaust all residue classes modulo $p$. Hence we may choose a residue class
\[
r_p \pmod p
\]
such that
\[
r_p\not\equiv i-jD\pmod p
\qquad (1\le i\le s,\ 0\le j\le t-1).
\]

By the Chinese remainder theorem, there exists an integer $a$ such that
\[
a\equiv r_p \pmod p
\qquad (\forall p\in\mathcal P),
\]
and also
\[
a>2(t-1)D+4s.
\]

For $1\le i\le s$ and $0\le j\le t-1$, define
\[
a_{i,j}:=a+jD-i.
\]
These integers are pairwise distinct. Indeed, if
\[
a+jD-i=a+\ell D-k,
\]
then
\[
(j-\ell)D=i-k.
\]
Since
\[
|i-k|<s\le st\le D,
\]
it follows that $j=\ell$, and then $i=k$.

Now set
\[
A:=\{a_{i,j}:1\le i\le s,\ 0\le j\le t-1\},
\qquad
N:=a+(t-1)D-1.
\]
Then $|A|=st$, and every element of $A$ lies in $\{1,\dots,N\}$, because
\[
1\le a-s\le a_{i,j}\le a+(t-1)D-1=N.
\]

We claim that for any two pairs $(i,j)$ and $(k,\ell)$,
\[
\gcd(a_{i,j},a_{k,\ell})\mid (k-i).
\]
Let
\[
d:=\gcd(a_{i,j},a_{k,\ell}).
\]
Since
\[
a_{i,j}-a_{k,\ell}=(j-\ell)D+(k-i),
\]
we have
\[
d\mid (j-\ell)D+(k-i).
\]

We distinguish three cases.

\medskip

\noindent
\emph{Case 1:} $i=k$.

Then
\[
d\mid 0=(k-i),
\]
so there is nothing to prove.

\medskip

\noindent
\emph{Case 2:} $i\ne k$ and $j=\ell$.

Then immediately
\[
d\mid (k-i).
\]

\medskip

\noindent
\emph{Case 3:} $i\ne k$ and $j\ne \ell$.

Let
\[
q:=|j-\ell|\in\{1,\dots,t-1\},
\]
and choose
\[
\delta=
\begin{cases}
k-i,& j>\ell,\\
i-k,& j<\ell.
\end{cases}
\]
Then
\[
d\mid qD+\delta,
\qquad
1\le |\delta|=|k-i|\le s-1.
\]

We first show that every prime divisor of $d$ is at most $st$. Suppose, for contradiction, that some prime $p>st$ divides $d$. Since
\[
p\mid d\mid qD+\delta,
\]
we have $p\in\mathcal P$ by definition of $\mathcal P$. On the other hand,
\[
p\mid d\mid a_{i,j}=a+jD-i,
\]
so
\[
a\equiv i-jD \pmod p,
\]
contradicting the choice of $a\equiv r_p\pmod p$, where $r_p$ avoids all residue classes $i-jD\pmod p$. Thus every prime divisor of $d$ is at most $st$.

Now let $p\le st$ be any prime. Since
\[
v_p(qD)\ge v_p(D)>v_p(\delta),
\]
it follows that
\[
v_p(qD+\delta)=v_p(\delta).
\]
Because $d\mid qD+\delta$, we obtain
\[
v_p(d)\le v_p(qD+\delta)=v_p(\delta)=v_p(k-i).
\]
Since all prime divisors of $d$ are at most $st$, we conclude that
\[
d\mid (k-i).
\]
This proves the claim.

Consequently, the system of congruences
\[
x\equiv -i \pmod{a_{i,j}}
\qquad (1\le i\le s,\ 0\le j\le t-1)
\]
is pairwise compatible: for any two congruences in the system, the difference of the residues is
\[
(-i)-(-k)=k-i,
\]
and by the claim this is divisible by $\gcd(a_{i,j},a_{k,\ell})$. Therefore, by the generalized Chinese remainder theorem~\cite[Theorem~11.5]{Donaldson180A}, the system has a solution $x_0$.

Replacing $x_0$ by a sufficiently large positive integer in the same residue class if necessary, we may assume that the interval introduced below lies in $[1,\infty)$.

Choose an integer $T$ such that
\[
2\bigl((t-1)D+s-1\bigr)<T\le a-2s.
\]
Such a choice is possible because
\[
a>2(t-1)D+4s.
\]

Consider the interval
\[
I:=(x_0-T,\ x_0-T+2N].
\]
Its length is $2N$.

Fix $(i,j)$ with $1\le i\le s$ and $0\le j\le t-1$. Since
\[
x_0\equiv -i\pmod{a_{i,j}},
\]
we have
\[
a_{i,j}\mid (x_0+i).
\]
Also, since
\[
a_{i,j}=a+jD-i,
\]
we have
\[
x_0+a+jD=(x_0+i)+a_{i,j},
\]
so $x_0+a+jD$ is another multiple of $a_{i,j}$.

We now show that these are exactly the two multiples of $a_{i,j}$ contained in $I$.

First, both of them lie in $I$. Indeed, the offset of $x_0+i$ from the left endpoint $x_0-T$ is
\[
T+i,
\]
which is positive and satisfies
\[
T+i\le T+s\le a-s<2N.
\]
Hence
\[
x_0+i\in I.
\]

Similarly, the offset of $x_0+a+jD$ from the left endpoint is
\[
T+a+jD,
\]
which is again positive. Moreover,
\[
T+a+jD\le T+a+(t-1)D\le (a-2s)+a+(t-1)D<2N,
\]
since
\[
2N=2a+2(t-1)D-2
\]
and
\[
(a-2s)+a+(t-1)D
<
2a+2(t-1)D-2
\]
for $s,t\ge 2$. Thus
\[
x_0+a+jD\in I.
\]

Next, the preceding multiple of $a_{i,j}$ is
\[
(x_0+i)-a_{i,j}.
\]
Its offset from the left endpoint is
\[
T+i-a_{i,j}.
\]
But
\[
a_{i,j}=a+jD-i\ge a-s,
\]
and
\[
T+i\le T+s\le a-s\le a_{i,j},
\]
so
\[
T+i-a_{i,j}\le 0.
\]
Hence this preceding multiple is not contained in $I$.

Finally, the next multiple after $x_0+a+jD$ is
\[
(x_0+i)+2a_{i,j},
\]
whose offset from the left endpoint is
\[
T+i+2a_{i,j}\ge T+1+2(a-s).
\]
Using
\[
T>2\bigl((t-1)D+s-1\bigr)=2(t-1)D+2s-2,
\]
we obtain
\[
T+i+2a_{i,j}
>
2(t-1)D+2s-2+1+2a-2s
=
2a+2(t-1)D-1
>
2N,
\]
because
\[
2N=2a+2(t-1)D-2.
\]
Therefore this multiple also lies outside $I$.

Thus $a_{i,j}$ has exactly two multiples in $I$, namely
\[
x_0+i
\qquad\text{and}\qquad
x_0+a+jD.
\]

It follows that every integer in $I$ divisible by some element of $A$ lies in the set
\[
B:=\{x_0+i:1\le i\le s\}\cup \{x_0+a+jD:0\le j\le t-1\}.
\]
The two parts of this union are disjoint, since
\[
x_0+i=x_0+a+jD
\]
would imply
\[
i=a+jD,
\]
which is impossible because $1\le i\le s<a\le a+jD$. Hence
\[
|B|=s+t.
\]

Therefore, inside the interval $I$, there are at most $s+t$ integers divisible by some element of $A$. In particular, there cannot exist a matching from $A$ into $I\cap\mathbb Z$ of size larger than $s+t$. Hence
\[
f(st)\le s+t.
\]

Now let
\[
n:=\lceil \sqrt m\rceil.
\]
Applying the above result with $s=t=n$, we obtain a set
\[
A\subseteq\{1,\dots,N\}
\]
of size $n^2$ and an interval $I$ of length $2N$ such that at most $2n$ integers in $I$ are divisible by some element of $A$.

Choose any subset $A'\subseteq A$ with $|A'|=m$. Then $A'\subseteq\{1,\dots,N\}$, and the set of integers in $I$ divisible by some element of $A'$ is contained in the corresponding set for $A$, so it still has cardinality at most $2n$. Therefore there cannot exist a matching from $A'$ into $I\cap\mathbb Z$ of size larger than $2n$, and hence
\[
f(m)\le 2n=2\lceil \sqrt m\rceil.
\]
This completes the proof.
\end{proof}





\begin{thebibliography}{9}

\bibitem{EP}
T. F. Bloom, Erd\H{o}s Problem \#650, \url{https://www.erdosproblems.com/650}, accessed 2026-03-06.

\bibitem{Donaldson180A}
N.~Donaldson,
\emph{Math 180A - Notes},
University of California, Irvine, March 14, 2018.
Available at:
\url{https://www.math.uci.edu/~ndonalds/math180a/notes.pdf}
(accessed March 7, 2026).


\bibitem{Er95c}
P. Erd\H{o}s, Some problems in number theory. Octogon Math. Mag. (1995), 3--5. 


\bibitem{ErSa59}
Erd\H{o}s, P. and Sar\'{a}nyi, Megjegyz\'{e}sek egy versenyfeladathoz. Matematikai Lapok (1959).



\end{thebibliography}









\end{document}